\documentclass[UTF8,12pt,fontset=windows]{ctexart}
\usepackage[paperwidth=240mm,paperheight=200mm,left=14mm,right=14mm,top=9mm,bottom=12mm]{geometry}
\usepackage{amsmath,amssymb,mathtools,bm,graphicx,xcolor}
\pagestyle{empty}
\setlength{\parindent}{0pt}
\setlength{\parskip}{4pt}
\setlength{\emergencystretch}{2em}
\xeCJKDeclareCharClass{CJK}{"2460 -> "24FF}
\xeCJKDeclareCharClass{CJK}{"2160 -> "2188}
\begin{document}
{\large\bfseries 2000.1}\quad {\small 原卷一（1）}

$\lim_{x \to 0} \dfrac{\arctan x - x}{\ln(1 + 2x^3)} = \_\_\_\_\_\_.$

\vfill
\newpage
{\large\bfseries 2000.2}\quad {\small 原卷一（2）}

设函数 $y = y(x)$ 由方程 $2^{xy} = x + y$ 所确定，则 $\left.dy\right|_{x=0} = \_\_\_\_\_\_.$

\vfill
\newpage
{\large\bfseries 2000.3}\quad {\small 原卷一（3）}

$\int_2^{+\infty} \dfrac{dx}{(x + 7)\sqrt{x - 2}} = \_\_\_\_\_\_.$

\vfill
\newpage
{\large\bfseries 2000.4}\quad {\small 原卷一（4）}

曲线 $y = (2x - 1)e^{\frac{1}{x}}$ 的斜渐近线方程为 $\_\_\_\_\_\_.$

\vfill
\newpage
{\large\bfseries 2000.5}\quad {\small 原卷一（5）}

设 $A = \begin{pmatrix} 1 & 0 & 0 & 0 \\ -2 & 3 & 0 & 0 \\ 0 & -4 & 5 & 0 \\ 0 & 0 & -6 & 7 \end{pmatrix}$，$E$ 为 4 阶单位矩阵，且 $B = (E + A)^{-1}(E - A)$，则 $(E + B)^{-1} = \_\_\_\_\_\_.$

\vfill
\newpage
{\large\bfseries 2000.6}\quad {\small 原卷二（1）}

设函数 $f(x) = \dfrac{x}{a + e^{bx}}$ 在 $(-\infty, +\infty)$ 内连续，且 $\lim_{x \to -\infty} f(x) = 0$，则常数 $a,b$ 满足（　）

(A) $a < 0, b < 0$.　　　　　(B) $a > 0, b > 0$.

(C) $a \leqslant 0, b > 0$.　　　　　(D) $a \geqslant 0, b < 0$.

\vfill
\newpage
{\large\bfseries 2000.7}\quad {\small 原卷二（2）}

设函数 $f(x)$ 满足关系式 $f''(x) + [f'(x)]^2 = x$，且 $f'(0) = 0$，则（　）

(A) $f(0)$ 是 $f(x)$ 的极大值.

(B) $f(0)$ 是 $f(x)$ 的极小值.

(C) 点 $(0, f(0))$ 是曲线 $y = f(x)$ 的拐点.

(D) $f(0)$ 不是 $f(x)$ 的极值，点 $(0, f(0))$ 也不是曲线 $y = f(x)$ 的拐点.

\vfill
\newpage
{\large\bfseries 2000.8}\quad {\small 原卷二（3）}

设函数 $f(x), g(x)$ 是大于零的可导函数，且 $f'(x)g(x) - f(x)g'(x) < 0$，则当 $a < x < b$ 时，有（　）

(A) $f(x)g(b) > f(b)g(x)$.　　　　　(B) $f(x)g(a) > f(a)g(x)$.

(C) $f(x)g(x) > f(b)g(b)$.　　　　　(D) $f(x)g(x) > f(a)g(a)$.

\vfill
\newpage
{\large\bfseries 2000.9}\quad {\small 原卷二（4）}

若 $\lim_{x \to 0} \dfrac{\sin 6x + xf(x)}{x^3} = 0$，则 $\lim_{x \to 0} \dfrac{6 + f(x)}{x^2}$ 为（　）

(A) 0.　　　　　(B) 6.　　　　　(C) 36.　　　　　(D) $\infty$.

\vfill
\newpage
{\large\bfseries 2000.10}\quad {\small 原卷二（5）}

具有特解 $y_1 = e^{-x}$，$y_2 = 2xe^{-x}$，$y_3 = 3e^x$ 的 3 阶常系数齐次线性微分方程是（　）

(A) $y''' - y'' - y' + y = 0$.

(B) $y''' + y'' - y' - y = 0$.

(C) $y''' - 6y'' + 11y' - 6y = 0$.

(D) $y''' - 2y'' - y' + 2y = 0$.

\vfill
\newpage
{\large\bfseries 2000.11}\quad {\small 原卷三}

设 $f(\ln x) = \dfrac{\ln(1 + x)}{x}$，计算 $\int f(x) \, dx$.

\vfill
\newpage
{\large\bfseries 2000.12}\quad {\small 原卷四}

设 $xOy$ 平面上有正方形 $D = \{(x,y) \mid 0 \leqslant x \leqslant 1,\ 0 \leqslant y \leqslant 1\}$ 及直线 $l: x + y = t \ (t \geqslant 0)$. 若 $S(t)$ 表示正方形 $D$ 位于直线 $l$ 左下方部分的面积，试求 $\int_0^x S(t) \, dt \ (x \geqslant 0)$.

\vfill
\newpage
{\large\bfseries 2000.13}\quad {\small 原卷五}

求函数 $f(x) = x^2 \ln(1 + x)$ 在 $x = 0$ 处的 $n$ 阶导数 $f^{(n)}(0) \ (n \geqslant 3)$.

\vfill
\newpage
{\large\bfseries 2000.14}\quad {\small 原卷六}

设函数 $S(x) = \int_0^x |\cos t| \, dt$，

(1) 当 $n$ 为正整数，且 $n\pi \leqslant x < (n + 1)\pi$ 时，证明：$2n \leqslant S(x) < 2(n + 1)$；

(2) 求 $\lim_{x \to +\infty} \dfrac{S(x)}{x}$.

\vfill
\newpage
{\large\bfseries 2000.15}\quad {\small 原卷七}

某湖泊的水量为 $V$，每年排入湖泊内含污染物 $A$ 的污水量为 $\dfrac{V}{6}$，流入湖泊内不含 $A$ 的水量为 $\dfrac{V}{6}$，流出湖泊的水量为 $\dfrac{V}{3}$. 已知 1999 年底湖中 $A$ 的含量为 $5m_0$，超过国家规定指标. 为了治理污染，从 2000 年初起，限定排入湖泊中含 $A$ 污水的浓度不超过 $\dfrac{m_0}{V}$. 问至多需经过多少年，湖泊中污染物 $A$ 的含量才可降至 $m_0$ 以内？（注：设湖泊中 $A$ 的浓度是均匀的.）

\vfill
\newpage
{\large\bfseries 2000.16}\quad {\small 原卷八}

设函数 $f(x)$ 在 $[0,\pi]$ 上连续，且 $\int_0^\pi f(x) \, dx = 0$，$\int_0^\pi f(x)\cos x \, dx = 0$. 试证明：在 $(0,\pi)$ 内至少存在两个不同的点 $\xi_1, \xi_2$，使 $f(\xi_1) = f(\xi_2) = 0$.

\vfill
\newpage
{\large\bfseries 2000.17}\quad {\small 原卷九}

已知 $f(x)$ 是周期为 5 的连续函数，它在 $x = 0$ 的某个邻域内满足关系式

$$f(1 + \sin x) - 3f(1 - \sin x) = 8x + \alpha(x),$$

其中 $\alpha(x)$ 是当 $x \to 0$ 时比 $x$ 高阶的无穷小，且 $f(x)$ 在 $x = 1$ 处可导，求曲线 $y = f(x)$ 在点 $(6, f(6))$ 处的切线方程.

\vfill
\newpage
{\large\bfseries 2000.18}\quad {\small 原卷十}

设曲线 $y = ax^2 \ (a > 0,\ x \geqslant 0)$ 与 $y = 1 - x^2$ 交于点 $A$，过坐标原点 $O$ 和点 $A$ 的直线与曲线 $y = ax^2$ 围成一平面图形. 问 $a$ 为何值时，该图形绕 $x$ 轴旋转一周所得的旋转体体积最大？最大体积是多少？

\vfill
\newpage
{\large\bfseries 2000.19}\quad {\small 原卷十一}

函数 $f(x)$ 在 $[0, +\infty)$ 上可导，$f(0) = 1$，且满足等式

$$f'(x) + f(x) - \frac{1}{x + 1}\int_0^x f(t) \, dt = 0.$$

(1) 求导数 $f'(x)$；

(2) 证明：当 $x \geqslant 0$ 时，不等式 $e^{-x} \leqslant f(x) \leqslant 1$ 成立.

\vfill
\newpage
{\large\bfseries 2000.20}\quad {\small 原卷十二}

设 $\boldsymbol{\alpha} = \begin{pmatrix} 1 \\ 2 \\ 1 \end{pmatrix}$，$\boldsymbol{\beta} = \begin{pmatrix} 1 \\ \dfrac{1}{2} \\ 0 \end{pmatrix}$，$\boldsymbol{\gamma} = \begin{pmatrix} 0 \\ 0 \\ 8 \end{pmatrix}$，$A = \boldsymbol{\alpha}\boldsymbol{\beta}^T$，$B = \boldsymbol{\beta}^T\boldsymbol{\alpha}$，其中 $\boldsymbol{\beta}^T$ 是 $\boldsymbol{\beta}$ 的转置，求解方程

$$2B^2 A^2 x = A^4 x + B^4 x + \boldsymbol{\gamma}.$$

\vfill
\newpage
{\large\bfseries 2000.21}\quad {\small 原卷十三}

已知向量组 $\boldsymbol{\beta}_1 = \begin{pmatrix} 0 \\ 1 \\ -1 \end{pmatrix}$，$\boldsymbol{\beta}_2 = \begin{pmatrix} a \\ 2 \\ 1 \end{pmatrix}$，$\boldsymbol{\beta}_3 = \begin{pmatrix} b \\ 1 \\ 0 \end{pmatrix}$ 与向量组 $\boldsymbol{\alpha}_1 = \begin{pmatrix} 1 \\ 2 \\ -3 \end{pmatrix}$，$\boldsymbol{\alpha}_2 = \begin{pmatrix} 3 \\ 0 \\ 1 \end{pmatrix}$，$\boldsymbol{\alpha}_3 = \begin{pmatrix} 9 \\ 6 \\ -7 \end{pmatrix}$ 具有相同的秩，且 $\boldsymbol{\beta}_3$ 可由 $\boldsymbol{\alpha}_1,\boldsymbol{\alpha}_2,\boldsymbol{\alpha}_3$ 线性表示，求 $a,b$ 的值.

\vfill
\end{document}
